\begin{longtable}{|p{0.276\textwidth}|p{0.092\textwidth}|p{0.46\textwidth}|p{0.092\textwidth}|}
\caption{Mandatory global attributes for GDS 2.0 netCDF data files}
\label{tab:global-attributes} \\
\hline \endfirsthead \endhead \endfoot \hline \endlastfoot
\rowcolor{blue!25} \textbf{Global Attribute Name} & \textbf{Type} & \textbf{Description} & \textbf{Source} \\ \hline
\rowcolor{cyan!25} 
acknowledgement & string & A place to acknowledge various types of support for the project that produced this data. & ACDD \\ \hline

\rowcolor{cyan!25} 
author & string & TBD & ECCO \\ \hline

\rowcolor{cyan!25} 
coordinates\_comment & string & TBD & ECCO \\ \hline

\rowcolor{cyan!25} 
cdm\_data\_type & string & The data type, as derived from Unidata's Common Data Model Scientific Data types and understood by THREDDS. (This is a THREDDS "dataType", and is different from the CF NetCDF attribute 'featureType', which indicates a Discrete Sampling Geometry file in CF.) & ACDD \\ \hline

\rowcolor{cyan!25} 
comment & string & Miscellaneous information about the data, not captured elsewhere. This attribute is defined in the CF Conventions. & CF, ACDD \\ \hline

\rowcolor{cyan!25} 
conventions & string & A comma-separated list of the conventions that the data follow. For files that follow this version of ACDD, include the string “ACDD-1.3” &  \\ \hline

\rowcolor{cyan!25} 
creator\_name & string & The name of the person (or other creator type specified by the creator\_type attribute) principally responsible for creating this data. & ACDD \\ \hline

\rowcolor{cyan!25} 
creator\_email & string & The email address of the person (or other creator type specified by the creator\_type attribute) principally responsible for creating this data. & ACDD \\ \hline

\rowcolor{cyan!25} 
creator\_type & string & Specifies the type of creator with one of the following: person, group, institution, or position. If this attribute is not specified, the creator is assumed to be a person & ACDD \\ \hline

\rowcolor{cyan!25} 
creator\_institution & string & The creator’s institution. The value should be specified even if it matches the value of publisher\_institution or if creator\_type is an institution & ACDD \\ \hline

\rowcolor{cyan!25} 
creator\_url & string & The URL of the of the person (or other creator type specified by the creator\_type attribute) principally responsible for creating this data. & ACDD \\ \hline

\rowcolor{cyan!25} 
date\_created & string & The date on which this version of the data was created. & ACDD \\ \hline

\rowcolor{cyan!25} 
date\_issued & string & The date on which this version of the data (including all modifications) was issued. & ACDD \\ \hline

\rowcolor{cyan!25} 
date\_modified & string & The date on which this version of the data was last modified. & ACDD \\ \hline

\rowcolor{cyan!25} 
date\_metadata\_modified & string & The date on which this version of the metadata was last modified. & ACDD \\ \hline

\rowcolor{cyan!25} 
geospatial\_lon\_max & float & Decimal degrees east, range -180 to +180.  This is equivalent to GDS “westernmost\_longitude” & ACDD \\ \hline

\rowcolor{cyan!25} 
geospatial\_lat\_resolution & float & Latitude Resolution in units matching geospatial\_lat\_units. & ACDD \\ \hline

\rowcolor{cyan!25} 
geospatial\_lat\_units & string & Units of the latitudinal resolution. Typically "degrees\_north" & ACDD \\ \hline

\rowcolor{cyan!25} 
geospatial\_lon\_resolution & float & Longitude Resolution in units matching geospatial\_lon\_resolution & ACDD \\ \hline

\rowcolor{cyan!25} 
geospatial\_lon\_units & string & Units of the longitudinal resolution. Typically "degrees\_east" & ACDD \\ \hline

\rowcolor{cyan!25} 
geospatial\_vertical\_min & string & Describes the numerically smaller vertical limit and can be part of a 2D or 3D bounding region.  See also geospatial\_vertical\_positive and geospatial\_vertical\_units & ACDD \\ \hline

\rowcolor{cyan!25} 
geospatial\_vertical\_max & string & Describes the numerically larger vertical limit and can be part of a 2D or 3D bounding region.  See also geospatial\_vertical\_positive and geospatial\_vertical\_units & ACDD \\ \hline

\rowcolor{cyan!25} 
geospatial\_vertical\_positive & string & Values include either “up” or “down.” If “up,” vertical values are interpreted as “altitude,” with negative values corresponding to below the reference datum (e.g., under water). If “down,” vertical values are interpreted as “depth,” positive values correspond to below the reference datum. Note that if geospatial\_vertical\_positive is “down” (depth orientation), geospatial\_vertical\_min specifies the data’s vertical location furthest from the Earth’s center, and geospatial\_vertical\_max specifies the location closest to the Earth’s center. & ACDD \\ \hline

\rowcolor{cyan!25} 
geospatial\_vertical\_resolution & string & Information about the targeted vertical spacing of points. Example: “25 meters” & ACDD \\ \hline

\rowcolor{cyan!25} 
geospatial\_vertical\_units & string & Units for the vertical axis described in geospatial\_vertical\_min and geospatial\_vertical\_max. The default is “EPSG:4979” (height above the ellipsoid, in meters), but other vertical coordinate reference systems may be specified. Note that the common oceanographic practice of using pressure for a vertical coordinate, while not strictly a depth, can be specified using the unit bar. Examples: “EPSG:5829” (instantaneous height above sea level), “EPSG:5831” (instantaneous depth below sea level) & ACDD \\ \hline

\rowcolor{cyan!25} 
history & string & The name of the institution principally responsible for originating this data. This attribute is recommended by the CF convention. & CF, ACDD \\ \hline

\rowcolor{cyan!25} 
id & string & An identifier for the data set, provided by and unique within its naming authority. The combination of the "naming authority" and the "id" should be globally unique, but the id can be globally unique by itself also. IDs can be URLs, URNs, DOIs, meaningful text strings, a local key, or any other unique string of characters. The id should not include white space characters. & ACDD \\ \hline

\rowcolor{cyan!25} 
institution & string & The name of the institution principally responsible for originating this data. This attribute is recommended by the CF convention. & CF, ACDD \\ \hline

\rowcolor{cyan!25} 
instrument\_vocabulary & string & Controlled vocabulary for the names used in instrument.  Example: “GCMD instrument keywords” & CF, ACDD \\ \hline

\rowcolor{cyan!25} 
keywords & string & GCMD Science Keyword(s) & ACDD \\ \hline

\rowcolor{cyan!25} 
keywords\_vocabulary & string & The unique name or identifier of the vocabulary from which keywords are taken. e.g., the NASA Global Change Master Directory (GCMD) Science Keywords. & ACDD \\ \hline

\rowcolor{cyan!25} 
license & string & Provide the URL to a standard or specific license, enter "Freely Distributed" or "None", or describe any restrictions to data access and distribution in free text. & ACDD \\ \hline

\rowcolor{cyan!25} 
platform & string & Name of the platform(s) that supported the sensor data used to create this data set or product. Platforms can be of any type, including satellite, ship, station, aircraft or other. Indicate controlled vocabulary used in platform\_vocabulary. & ACDD \\ \hline

\rowcolor{cyan!25} 
platform\_vocabulary & string & Controlled vocabulary for the names used in the "platform" attribute. & ACDD \\ \hline

\rowcolor{cyan!25} 
metadata\_link & string & Link to collection metadata record at archive & ACDD \\ \hline

\rowcolor{cyan!25} 
naming\_authority & string & The organization that provides the initial id (see above) for the dataset. The naming authority should be uniquely specified by this attribute via reverse-DNS naming convention. & ACDD \\ \hline

\rowcolor{cyan!25} 
netcdf\_version\_id  & string & Version of netCDF libraries used to create this file. For example, "4.1.1" & GDS \\ \hline

\rowcolor{cyan!25} 
geospatial\_bounds\_crs & string & The CRS of the point coordinates in geospatial\_bounds. & GDS \\ \hline

\rowcolor{cyan!25} 
geospatial\_lat\_max & float & Decimal degrees north, range -90 to +90.  This is equivalent to GDS “northernmost latitude” & ACDD \\ \hline

\rowcolor{cyan!25} 
processing\_level & string & A textual description of the processing (or quality control) level of the data. & ACDD \& GDS \\ \hline

\rowcolor{cyan!25} 
product\_name & string & TBD & ECCO \\ \hline

\rowcolor{cyan!25} 
product\_time\_coverage\_end & string & TBD & ECCO \\ \hline

\rowcolor{cyan!25} 
product\_time\_coverage\_start & string & TBD & ECCO \\ \hline

\rowcolor{cyan!25} 
product\_version & string & The product version of this data file & GDS \\ \hline

\rowcolor{cyan!25} 
program & string & The overarching program of which the data belongs. A program consists of a set (or portfolio) of related and possibly interdependent projects that meet an overarching objective.  Example: 'NASA Physical Oceanography, Cryosphere, Modeling, Analysis, and Prediction (MAP)' & GDS \\ \hline

\rowcolor{cyan!25} 
project & string & The name of the project(s) principally responsible for originating this data.  Example: 'Estimating the Circulation and Climate of the Ocean (ECCO)' & ACDD \\ \hline

\rowcolor{cyan!25} 
publisher\_email & string & The email address of the person (or other entity specified by the publisher\_type attribute) responsible for publishing the data file or product to users, with its current metadata and format. & ACDD \\ \hline

\rowcolor{cyan!25} 
publisher\_institution & string & The publisher’s institution. If publisher\_type is an institution, this attribute should have the same value as publisher\_name.  Example: “PO.DAAC” & ACDD \\ \hline

\rowcolor{cyan!25} 
publisher\_name & string & The name of the person (or other entity specified by the publisher\_type attribute) responsible for publishing the data file or product to users, with its current metadata and format. & ACDD \\ \hline

\rowcolor{cyan!25} 
publisher\_type & string & Specifies type of publisher with one of the following: person, group, institution, or position. If this attribute is not specified, the publisher is assumed to be a person.  Example: “institution” & ACDD \\ \hline

\rowcolor{cyan!25} 
publisher\_url & string & The URL of the person (or other entity specified by the publisher\_type attribute) responsible for publishing the data file or product to users, with its current metadata and format. & ACDD \\ \hline

\rowcolor{cyan!25} 
references & string & Published or web-based references that describe the data or methods used to produce it. Recommend URIs (such as a URL or DOI) for papers or other references. This attribute is defined in the CF conventions. & ACDD \\ \hline

\rowcolor{cyan!25} 
source & string & Method of production of the original data. & CF \\ \hline

\rowcolor{cyan!25} 
geospatial\_lat\_min & float & Decimal degrees north, range -90 to +90.  This is equivalent to GDS “southernmost\_latitude” & ACDD \\ \hline

\rowcolor{cyan!25} 
spatial\_resolution & string & A string describing the approximate resolution of the product. & GDS \\ \hline

\rowcolor{cyan!25} 
standard\_name\_vocabulary & string & The name and version of the controlled vocabulary from which variable standard names are taken. & ACDD \\ \hline

\rowcolor{cyan!25} 
start\_time & string & Representative date and time of the end of the granule in the ISO 8601 compliant format of "yyyymmddThhmmssZ". & GDS \\ \hline

\rowcolor{cyan!25} 
stop\_time & string & Representative date and time of the end of the granule in the ISO 8601 compliant format of "yyyymmddThhmmssZ". & GDS \\ \hline

\rowcolor{cyan!25} 
summary & string & A paragraph describing the dataset, analogous to an abstract for a paper. & ACDD \\ \hline

\rowcolor{cyan!25} 
time\_coverage\_end & string & Describes the time of the last data point in the data. Use ISO 8601:2004 date format, preferably the extended format as recommended in ACDD Section 2.6.  Example: “2017-12-30T00:00:00” & ACDD \\ \hline

\rowcolor{cyan!25} 
time\_coverage\_start & string & Describes the time of the first data point in the data. Use the ISO 8601:2004 date format, preferably the extended format as recommended in ACDD Section 2.6.  Example: “2017-12-29T00:00:00” & ACDD \\ \hline

\rowcolor{cyan!25} 
time\_coverage\_duration & string & Describes the duration of the data. Use ISO 8601:2004 duration format, preferably the extended format as recommended in ACDD Section 2.6.  Example: “P1D” & ACDD \\ \hline

\rowcolor{cyan!25} 
time\_coverage\_resolution & string & Describes the targeted time period between each value in the data. Use ISO 8601:2004 duration format, preferably the extended format as recommended in ACDD Section 2.6.  Example: “P1D” & ACDD \\ \hline

\rowcolor{cyan!25} 
title & string & A short phrase or sentence describing the dataset. In many discovery systems, the title will be displayed in the results list from a search, and therefore should be human readable and reasonable to display in a list of such names. This attribute is recommended by the NetCDF Users Guide (NUG) and the CF conventions. & CF, ACDD \\ \hline

\rowcolor{cyan!25} 
uuid & string & A Universally Unique Identifier (UUID). Numerous, simple tools can be used to create a UUID, which is inserted as the value of this attribute. See http://en.wikipedia.org/wiki/Universally\_Unique\_Identifier for more information and tools. & GDS \\ \hline

\rowcolor{cyan!25} 
geospatial\_lon\_min & float & Decimal degrees east, range -180 to +180.  This is equivalent to GDS “easternmost\_longitude” & ACDD \\ \hline

\end{longtable}
